\documentclass[11pt]{article}
\usepackage[margin=1in]{geometry}
\usepackage{amsmath,amssymb,amsthm,mathtools}
\usepackage{graphicx}
\usepackage{microtype}
\usepackage[colorlinks=true,linkcolor=blue,urlcolor=blue,citecolor=blue]{hyperref}

\newtheorem{theorem}{Theorem}
\newtheorem{lemma}[theorem]{Lemma}
\newcommand{\R}{\mathbb R}
\newcommand{\1}{\mathbf 1}
\newcommand{\Mor}{L^{p,\lambda}(\R^n)}
\newcommand{\Vast}{V^{(\ast)}L^{p,\lambda}(\R^n)}

\title{Singular-type operators preserve the Morrey vanishing property $(V^\ast)$}
\author{Candidate full answer to the question in arXiv:1911.06551}
\date{17 August 2026}

\begin{document}
\maketitle

\noindent\fbox{\parbox{0.94\linewidth}{\textbf{Status: candidate full proof, likely valid.}
For every $1<p<\infty$ and $0\leq\lambda<n$, every $p$-admissible
sublinear singular-type operator preserves $V^{(\ast)}L^{p,\lambda}(\R^n)$.
This includes all singular integral operators asked about in the source.}}

\section{Source question and definitions}

For $1<p<\infty$ and $0\leq\lambda<n$, the Morrey space $\Mor$ has norm
\[
 \|f\|_{p,\lambda}
 :=\sup_{x\in\R^n,\ r>0}
 \left(r^{-\lambda}\int_{B(x,r)}|f(y)|^p\,dy\right)^{1/p}.
\]
Writing $\chi_N=\1_{\R^n\setminus B(0,N)}$, define
\[
 \mathcal A_{N,p}(f)
 :=\sup_{x\in\R^n}\int_{B(x,1)}|f(y)|^p\chi_N(y)\,dy.
\]
The subspace $\Vast$ consists of the $f\in\Mor$ satisfying
$\mathcal A_{N,p}(f)\to0$.

The source calls a sublinear operator $T$ \emph{$p$-admissible singular-type}
if it is bounded on $L^p(\R^n)$ and, for $z\in B(x,r)$,
\begin{equation}\label{eq:size}
 \left|T\!\left(f\1_{\R^n\setminus B(x,2r)}\right)(z)\right|
 \leq C_T\int_{\R^n\setminus B(x,2r)}
       \frac{|f(y)|}{|y-z|^n}\,dy .
\end{equation}
Its Lemma 2.3 implies that such $T$ is bounded on $\Mor$.
Immediately before Theorem 3.11, the authors ask whether a singular integral
$S$ in this class preserves $(V^\ast)$ \cite[p.~9]{source}.

\begin{figure}[ht]
\centering
\includegraphics[width=\linewidth]{figures/open_problem_crop.png}
\caption{The exact open question on page 9 of the arXiv PDF.}
\end{figure}

\section{Proof intuition}

The size condition reduces the only difficult piece to an order-$n$ tail
potential.  On a dyadic annulus of radius $2^k$, the $(V^\ast)$ hypothesis
gives a bound by the small far-unit-ball quantity $\eta_N$, but this bound
alone is not summable in $k$.  The Morrey norm gives the complementary bound
$M2^{-k(n-\lambda)/p}$.  Taking the better estimate on every annulus and
summing gives
\[
 \sum_{k\geq0}\min\{\eta_N,M2^{-k(n-\lambda)/p}\}
 \lesssim \eta_N\bigl(1+\log(M/\eta_N)\bigr)\longrightarrow0.
\]
This logarithmic interpolation is the missing critical-tail mechanism.

\section{Full positive answer}

\begin{lemma}[Critical tail estimate]\label{lem:tail}
Let $f\in\Mor$, $M=\|f\|_{p,\lambda}$, and
\[
 F_N=f\1_{\R^n\setminus B(0,N)},\qquad
 \eta_N=\sup_{u\in\R^n}\|F_N\|_{L^p(B(u,1))}.
\]
Set $a=(n-\lambda)/p>0$.  If $G$ is any measurable function with
$|G|\leq|F_N|$ and $G=0$ on $B(x,2)$, then, uniformly for $z\in B(x,1)$,
\begin{equation}\label{eq:tail-bound}
 \int_{\R^n}\frac{|G(y)|}{|y-z|^n}\,dy
 \leq C_{n,p,\lambda}\,eta_N
 \left(1+\log_+\frac{M}{\eta_N}\right),
\end{equation}
with the right side interpreted as $0$ when $\eta_N=0$.
\end{lemma}

\begin{proof}
Every ball of radius $R\geq1$ can be covered by at most $C_nR^n$ unit
balls.  Since multiplication by an indicator does not increase the Morrey
norm,
\begin{equation}\label{eq:two-ball-bounds}
 \|G\|_{L^p(B(z,R))}
 \leq C_n\min\{\eta_NR^{n/p},\,MR^{\lambda/p}\}.
\end{equation}
The support condition and $z\in B(x,1)$ imply $|y-z|\geq1$ on the support of
$G$.  Split that support into the annuli
$2^k\leq|y-z|<2^{k+1}$, $k\geq0$.  H\"older's inequality and
\eqref{eq:two-ball-bounds} yield
\[
 \int_{2^k\leq|y-z|<2^{k+1}}
       \frac{|G(y)|}{|y-z|^n}\,dy
 \leq C_{n,p}\min\{\eta_N,\,M2^{-ak}\}.
\]
Here the first alternative follows because the volume factor
$2^{kn/p'}$ and the unit-cover factor $2^{kn/p}$ exactly cancel the kernel
$2^{-kn}$; the second leaves $2^{-k(n-\lambda)/p}$.

If $0<\eta_N\leq M$, let
\[
 K=\left\lceil\frac{\log(M/\eta_N)}{a\log2}\right\rceil.
\]
The terms with $k<K$ total at most $K\eta_N$, while the remaining geometric
tail is at most
$M2^{-aK}/(1-2^{-a})\leq\eta_N/(1-2^{-a})$.
This proves \eqref{eq:tail-bound}.  The cases $M=0$ or $\eta_N=0$ are
immediate.  (In fact $\eta_N\leq M$ by the Morrey estimate at radius one.)
\end{proof}

\begin{theorem}\label{thm:main}
Let $1<p<\infty$ and $0\leq\lambda<n$.  Every $p$-admissible sublinear
singular-type operator $T$ maps $\Vast$ boundedly into itself.  In
particular, every singular integral operator $S$ considered in the source
preserves the vanishing property $(V^\ast)$.
\end{theorem}

\begin{proof}
The source's Morrey boundedness result already gives $Tf\in\Mor$ and
$\|Tf\|_{p,\lambda}\lesssim\|f\|_{p,\lambda}$.  It remains to prove
$(V^\ast)$.

Fix $f\in\Vast$, write $M=\|f\|_{p,\lambda}$, and use $\eta_N$ from
Lemma~\ref{lem:tail}.  Thus $\eta_N=\mathcal A_{N,p}(f)^{1/p}\to0$.
For fixed $N$ and $|x|>N+3$, decompose
\[
 f=u_x+v_N+w_{N,x},
\]
where
\[
 u_x=f\1_{B(x,2)},\qquad
 v_N=f\1_{B(0,N)},\qquad
 w_{N,x}=f\1_{(\R^n\setminus B(0,N))\setminus B(x,2)}.
\]
The first two supporting balls are disjoint, so these three pieces cover
$f$.  By sublinearity and Minkowski's inequality it suffices to estimate the
three outputs on $B(x,1)$.

First, the global $L^p$ boundedness of $T$ and a fixed finite cover of
$B(x,2)$ by unit balls give
\begin{equation}\label{eq:local}
 \|Tu_x\|_{L^p(B(x,1))}
 \leq C_T\|u_x\|_{L^p(\R^n)}
 \leq C_{n,T}\eta_N,
\end{equation}
because $B(x,2)\subset\R^n\setminus B(0,N)$.

Second, $v_N$ is supported off $B(x,2)$.  For $z\in B(x,1)$, the size
condition gives
\[
 |Tv_N(z)|
 \leq C_T\int_{B(0,N)}\frac{|f(y)|}{|y-z|^n}\,dy
 \leq \frac{C_T\|f\1_{B(0,N)}\|_{L^1}}
              {( |x|-N-1)^n}.
\]
Consequently, for each fixed $N$,
\begin{equation}\label{eq:origin}
 \|Tv_N\|_{L^p(B(x,1))}\longrightarrow0
 \qquad (|x|\to\infty).
\end{equation}

Finally, $w_{N,x}$ satisfies the hypotheses of Lemma~\ref{lem:tail}; applying
the size condition and then the lemma gives
\begin{equation}\label{eq:tailout}
 \|Tw_{N,x}\|_{L^p(B(x,1))}
 \leq C_{n,p,\lambda,T}\eta_N
 \left(1+\log_+\frac{M}{\eta_N}\right).
\end{equation}
From \eqref{eq:local}--\eqref{eq:tailout},
\[
 \limsup_{|x|\to\infty}\|Tf\|_{L^p(B(x,1))}
 \leq C\eta_N\left(1+\log_+\frac{M}{\eta_N}\right).
\]
Letting $N\to\infty$ makes the right side zero.  Hence the $L^p$ norm of
$Tf$ on unit balls tends to zero uniformly as their centers go to infinity.

To recover the source's cut-off formulation, note that
$B(x,1)\cap(\R^n\setminus B(0,R))\neq\varnothing$ implies $|x|>R-1$.
Therefore
\[
 \mathcal A_{R,p}(Tf)^{1/p}
 \leq\sup_{|x|>R-1}\|Tf\|_{L^p(B(x,1))}\longrightarrow0.
\]
Thus $Tf\in\Vast$.
\end{proof}

\section{Scope, verification, and novelty}

The theorem uses exactly the source definition of a $p$-admissible
singular-type operator.  It does not require linearity, kernel smoothness, or
additional cancellation.  The strict inequality $\lambda<n$ is essential to
this argument because it makes $a=(n-\lambda)/p$ positive.  All source
parameters, including $\lambda=0$, are covered.

The dyadic transition and vanishing modulus were stress-tested over 72,000
parameter triples by the included verifier; this is a check of the algebra,
not a proof.  The formal proof above supplies the exact geometric-series
argument.

A bounded novelty search on 17 August 2026 covered the local run indexes,
the arXiv identifier and title, the exact phrase ``we do not know whether
$S$ preserves the vanishing property $(V^\ast)$'', the authors, singular
operators on vanishing Morrey spaces, and later citing papers.  The closest
item is the same authors' generalized-Morrey follow-up \cite{followup}.
Crucially, its Section 5 again states the same uncertainty and its Theorem 5.3
only proves invariance of the closure of $C_0^\infty$, not invariance of
$V^{(\ast)}$ itself.  Searches found no later full answer or matching
logarithmic dyadic-tail argument.  Novelty confidence is moderate rather than
high because the argument is elementary and the search was bounded.

\paragraph{Human-review recommendation.}
Check (i) that the source size condition applies to the noncompact truncated
tail $w_{N,x}$ under its standard Morrey extension, (ii) the annular H\"older
exponents in Lemma~\ref{lem:tail}, and (iii) the last equivalence between
far-center decay and the cut-off definition of $(V^\ast)$.

\begin{thebibliography}{2}
\bibitem{source}
A.~\c{C}. Alabalik, A.~Almeida, and S.~Samko,
\emph{On the invariance of certain vanishing subspaces of Morrey spaces with
respect to some classical operators}, Banach J. Math. Anal. 14 (2020),
987--1000; arXiv:1911.06551,
\url{https://arxiv.org/abs/1911.06551}.

\bibitem{followup}
A.~\c{C}. Alabalik, A.~Almeida, and S.~Samko,
\emph{Preservation of certain vanishing properties of generalized Morrey
spaces by some classical operators}, Math. Methods Appl. Sci. 43 (2020),
9375--9386, \href{https://doi.org/10.1002/mma.6235}{doi:10.1002/mma.6235}.
\end{thebibliography}

\end{document}
